\chapter{开发环境、编译器诊断和学习实验}

\section{问题入口：同一段代码为什么有人能编译，有人不能}

学习 \cpp 时，经常遇到“我这里能编译，你那里不行”。原因可能不是代码本身，而是编译器版本、标准选项、标准库实现、构建类型和警告设置不同。下面这段代码使用了 \cpp20 的 \code{std::span}：

\begin{lstlisting}[language=C++,caption={依赖 Cpp20 的代码}]
#include <span>
#include <vector>

int sum(std::span<const int> values)
{
    int total = 0;
    for (int value : values) {
        total += value;
    }
    return total;
}
\end{lstlisting}

如果编译命令没有写 \code{-std=c++20}，编译器可能根本不认识 \code{std::span}。因此从第一天开始，构建命令就应该是学习记录的一部分，而不是“随便点一下运行”。

\section{固定最小编译命令}

单文件实验建议从下面命令开始：

\begin{lstlisting}[language=bash,caption={推荐的单文件实验命令}]
g++ -std=c++20 -Wall -Wextra -Wpedantic -Wconversion main.cpp -o app
./app
\end{lstlisting}

\code{-std=c++20} 固定语言版本。\code{-Wall -Wextra -Wpedantic} 打开常见警告。\code{-Wconversion} 会提示一些隐式转换风险。不要把警告当成“编译器太烦”。它们经常指出真实问题：

\begin{lstlisting}[language=C++,caption={隐式转换可能丢信息}]
double average = 95.8;
int score = average;
\end{lstlisting}

如果业务上就是要截断，应该显式写出转换，并在变量名或附近逻辑中说明为什么可以截断。隐式发生的丢失信息，最容易藏 bug。

\section{读错误信息的顺序}

编译错误信息可能很长，尤其包含模板时。不要从第一行开始逐字翻译，先按这个顺序找：

\begin{enumerate}
  \item 第一个出现你自己文件路径的位置。
  \item 报错行附近的真正表达式。
  \item 错误类别：类型不匹配、找不到名字、没有匹配函数、约束不满足。
  \item 后续 note 里给出的候选函数或约束原因。
\end{enumerate}

例如：

\begin{lstlisting}[language=C++,caption={类型不匹配示例}]
std::vector<int> values{1, 2, 3};
std::string text = values;
\end{lstlisting}

编译器可能打印一串构造函数候选，但核心错误很简单：\code{std::string} 没有办法从 \code{std::vector<int>} 构造。先抓核心类别，再看候选细节。

\section{Debug 和 Release 的区别}

Debug 构建更适合调试，Release 构建更适合性能测量。不要用 Debug 时间判断算法快慢：

\begin{lstlisting}[language=bash,caption={两种构建模式}]
g++ -std=c++20 -O0 -g main.cpp -o app-debug
g++ -std=c++20 -O2 -DNDEBUG main.cpp -o app-release
\end{lstlisting}

\code{-O0 -g} 保留调试信息，不做太多优化。\code{-O2} 会做优化。\code{-DNDEBUG} 会关闭 \code{assert}。这推出一个重要规则：\code{assert} 只能检查内部不变量，不能替代用户输入校验。因为 Release 下它可能根本不执行。

\section{用小实验验证规则}

学 \cpp 不应该只背结论。每个规则都可以设计小实验。例如验证 \code{vector} 扩容导致指针失效：

\begin{lstlisting}[language=C++,caption={观察 vector 扩容}]
std::vector<int> values;
values.reserve(1);
values.push_back(10);
const int* old_address = values.data();
values.push_back(20);
const int* new_address = values.data();
std::cout << (old_address == new_address) << '\n';
\end{lstlisting}

这段代码不保证一定输出 \code{0}，因为实现可以有不同扩容策略。但当输出 \code{0} 时，你就能直观看到原指针不再指向同一块存储。实验不是替代标准，而是帮助你把抽象规则和运行现象连起来。

\section{最小项目脚手架}

当单文件开始变长，就应拆成最小项目：

\begin{lstlisting}[caption={最小项目目录}]
demo/
  CMakeLists.txt
  include/demo/math.hpp
  src/math.cpp
  tests/math_tests.cpp
\end{lstlisting}

构建命令固定为：

\begin{lstlisting}[language=bash,caption={固定项目构建命令}]
cmake -S . -B build -DCMAKE_BUILD_TYPE=Debug
cmake --build build
ctest --test-dir build --output-on-failure
\end{lstlisting}

这三个命令让项目能被别人重复构建。能重复构建，问题才容易复现；能复现，才谈得上调试和修复。

\begin{exercisebox}
建立一个 \filepath{demo} 项目，实现 \code{sum(std::span<const int>)}。先故意不写 \code{-std=c++20}，观察错误；再加上标准选项。然后加入一个会触发 \code{-Wconversion} 的例子，写下警告真正提醒了什么风险。
\end{exercisebox}
